\section{개요}
\label{sec: start}
Problem Solving에 필요한 알고리즘 공부에 앞서, 알고리즘을 평가하는 방법을 알 필요가 있다. 보통 알고리즘의 평가는 \textbf{시간과 메모리}의 두 가지 측면에서 이루어진다. Problem Solving에서 시간 제한과 메모리 제한은 문제를 풀기 전에 꼭 확인해야 하는 요소 중 하나이며, 시간 제한에 따라 문제 풀이가 달라지기도 한다.

당연하게도 시간과 메모리는 제약 조건(Constraint)에 대한 함수로 표시될 것이다. 아주 간단히 1부터 $n$까지 수의 합을 구하는 문제와 이에 대한 아래 두 가지 알고리즘을 보자:

\begin{itemize}
    \item 알고리즘 1: 반복문을 이용해 $1$부터 $n$까지 순서대로 구한다.
    \item 알고리즘 2: $\sum_{i=1}^n i = \frac{n(n+1)}{2}$ 공식을 이용해 구한다.
\end{itemize}

직관적으로 알고리즘 1보다 알고리즘 2가 더 빠른 시간에 작동할 것이라고 추측할 수 있고, 실제로 그렇다(궁금한 사람은 $n=100,000,000$ 정도를 대입하여 직접 실행해보자). 그렇다면 정학히 얼마나 더 빠를까?

이에 대한 답을 명확히 줄 수 있는 수학적 기호가 바로 시간 복잡도와 공간 복잡도이다. 알고리즘 1의 시간 복잡도와 공간 복잡도를 보면 반복문에서 $n$번의 연산을 하고, 답에 저장하는 변수 하나만 필요하므로 메모리는 $1$이 필요하다고 생각할 수 있다. 반면 알고리즘 2의 경우에는 반복문 없이 총 세 번의 연산으로 답을 구할 수 있다. 이 알고리즘의 경우에는 시간 복잡도는 $3$로 줄어들고 공간 복잡도는 여전히 $1$이다. 즉, 알고리즘 2가 알고리즘 1보다 $n$배 빠르다고 말할 수 있다.

이런 의문이 들었을 수도 있다: 정확히는 $n/3$배 빠른 것 아닌가? 위에서 $n$으로 표시한 이유는 시간 복잡도와 공간 복잡도에서 상수가 크게 중요하지 않기 때문이다. 시간 복잡도와 공간 복잡도는 정확히 얼마의 시간이 걸리는지 알기 위한 것이 아니라, 제약조건에 따라 어느 정도로 시간이 빠르게 증가하는지를 확인하고자 하는 것이다. 즉, 우리는 시간에 가장 큰 영향을 주는 항만 남겨 표기해야 한다. 그리고 이를 수학적으로 나타내기 위해 점근 표기법(asymptotic notation)을 사용한다.

\section{점근 표기법}
점근 표기법이란, 함수가 증가 또는 감소하는 경향성을 나타내기 위한 표기법이다. 대표적으로는 Big-O, small-o, Big-Theta, small-theta, Big-Omega, small-omega 표기법이 있으며, 이 중 이 책에서는 PS에서 자주 쓰이는 Big-O와 Big-Theta notation에 대해서만 서술할 것이다. 이들의 정의는 다음과 같다.
\begin{definition}[Big-O notation]
    $\limsup_{x\to\infty} f(x)/g(x)<\infty$이면 $f(x)=O(g(x))$이다.
\end{definition}
\begin{definition}[Big-Theta notation]\label{def: big-theta}
    $\limsup_{x\to\infty} f(x)/g(x)<\infty$이고 $\liminf_{x\to\infty} f(x)/g(x)$ $>0$이면 $f(x)=\Theta(g(x))$이다.
\end{definition}
위 정의에서 보다시피 Big-O notation은 $f(x)$의 점근적 상한, Big-Theta notation은 $f(x)$의 점근적 상한이자 하한을 나타내는 것으로, Big-Theta가 더 엄밀한 값을 나타낸다는 것을 알 수 있다. 하지만 PS에서는 이 둘을 혼용해서 모두 \Cref{def: big-theta}의 정의로 사용하며, 이 책에서도 이후로 나오는 Big-O notation에 대한 정의는 예의 정의로 쓰일 것이다.

이제 시간 복잡도를 표기하는 방법을 알아보기 위해 앞의 예시를 다시 한 번 가져와보자.
\begin{itemize}
    \item 알고리즘 1: 반복문을 이용해 $1$부터 $n$까지 순서대로 구한다.
    \item 알고리즘 2: $\sum_{i=1}^n i = \frac{n(n+1)}{2}$ 공식을 이용해 구한다.
\end{itemize}
\cref{sec: start}에서 분석한 내용과 같이 알고리즘 1에는 $n$번의 연산이 쓰이며, 알고리즘 2에는 3번의 연산이 쓰인다. 이를 Big-O notation으로 표현한다면 알고리즘 1은 $O(n)$의 시간 복잡도를, 알고리즘 2는 $O(1)$의 시간 복잡도를 가진다고 표현할 수 있다(\Cref{def: big-theta}에서는 계수가 포함되어도 상관이 없지만, 이제부터는 그 상수를 생략하고 표기할 것이다).

공간 복잡도에 대해서도 같은 표기법을 사용한다는 것을 알아두자.

\section{제한 시간과 시간 복잡도}
현재 컴퓨터는 1초에 약 $10^8$번 정도, 많게는 $10^9$번 정도의 연산을 수행할 수 있다고 알려져 있다. PS 문제를 푼다면 본인이 생각한 풀이의 시간 복잡도를 분석하여 그 식에 입력 제한을 대입하여 제한 시간 안에 풀이가 답을 찾을지 확인하는 것이 좋다. 다음은 PS에서 사용되는 대표적인 시간 복잡도이다.
\begin{itemize}
    \item 상수 시간 복잡도, $O(1)$
    \item 로그 시간 복잡도, $O(\log N), O(\log^2 N)$ 등
    \item 루트 시간 복잡도, $O(\sqrt N)$
    \item 다항 시간 복잡도, $O(N), O(N^2)$ 등
    \item 지수 시간 복잡도, $O(2^N), O(3^N)$ 등
    \item 팩토리얼 시간 복잡도, $O(N!)$
\end{itemize}
보통 시간 복잡도는 위 예시들의 합 또는 곱으로 이루어진다.

\section{연습문제}
\begin{problem}
    크기 $N$의 어떤 집합 $S$의 모든 부분집합을 출력하려고 한다. 완전 탐색 방식의 시간복잡도는 무엇인가?
\end{problem}
\begin{problem}
    $N$개의 노드로 이루어진 완전 이진 트리의 리프에서 루트까지 올라가는 경로를 출력하고자 한다. 시간복잡도는 무엇인가?
\end{problem}

(추가될 예정)